% !TeX program = xelatex
\documentclass[tikz,12pt]{standalone}
\usepackage[UTF8,scheme=plain,fontset=windows]{ctex}
\usepackage{tikz}
\usetikzlibrary{positioning,arrows.meta,calc}
\begin{document}
\begin{tikzpicture}[>=Latex,node distance=8mm and 10mm]
\tikzstyle{box}=[draw,rounded corners,align=center,minimum height=9mm,minimum width=28mm,font=\small]
\tikzstyle{widebox}=[draw,rounded corners,align=center,minimum height=9mm,minimum width=44mm,font=\small]
\node[box] (app1) {量子模拟\\量子化学};
\node[box,right=of app1] (app2) {优化\\采样/决策};
\node[box,right=of app2] (app3) {机器学习\\数据分析};

\node[widebox,below=10mm of app2] (arch) {体系结构转向\\单机扩展 $\rightarrow$ 模块化/分布式};

\node[box,below left=10mm and 2mm of arch] (node1) {本地处理模块\\高保真门/纠错};
\node[box,below=10mm of arch] (node2) {网络互连\\远程纠缠资源};
\node[box,below right=10mm and 2mm of arch] (node3) {量子接口\\物质 $\leftrightarrow$ 光子};

\draw[->,thick] (app1.south) -- (arch.north);
\draw[->,thick] (app2.south) -- (arch.north);
\draw[->,thick] (app3.south) -- (arch.north);

\draw[->,thick] (arch.south) -- (node1.north);
\draw[->,thick] (arch.south) -- (node2.north);
\draw[->,thick] (arch.south) -- (node3.north);

\end{tikzpicture}
\end{document}
